


Generated by \texttt{python -m crashguard.rd\_x13 analyse}. Design, scenario S-X and the agreement rule: \texttt{PREREGISTRATION\_V12.md} addendum B section 10.4. Constants: \texttt{ASSUMPTIONS.md}. Build and local ns-3 changes: \texttt{XSIM\_BUILD.md}. The calibration criterion of X1.2 was NOT MET (\texttt{XSIM\_CALIBRATION.md}), so every statement below is made under a calibration that did not reach its own criterion.

Scenario: the 32 SF9 cells of FRESH-A (SF11 is not expressible in ns-3's ALOHA region), 100 simulated hours per cell and policy, new seeds, Critical events only, frozen D1 table and parameters (nothing re-tuned). These cells were used once before for the custom-simulator verdicts; this run is a cross-check, not a second held-out test. ns-3 runs that crashed inside the module (a segmentation fault after an ACK callback; about 0.2\% of runs, same arguments crash again; cause not found) are dropped for all policies of that cell and hour: 22 of 3200 hours.

\subsection*{Verdicts of the V12 criteria on S-X}

{\scriptsize\raggedright
\begin{longtable}{p{0.237\linewidth}p{0.237\linewidth}p{0.237\linewidth}p{0.237\linewidth}}
\toprule
Hypothesis & Custom simulator & ns-3 & Same verdict \\
\midrule
\endhead
H-D1a no cell below the severity table by more than 1 point & met & not met & \textbf{no: simulator-dependent} \\
H-D1b at least 5 points above the severity table at loss of 40\% or more & met & not met & \textbf{no: simulator-dependent} \\
H-D1c duty limit never exceeded & met & met & yes \\
H-D1s not below the best baseline of the cell by more than 2 points & not met & not met & yes \\
\bottomrule
\end{longtable}}

A hypothesis whose verdict differs between the simulators is simulator-dependent and cannot be in the abstract. D1 enters the main paper only if H-D1a and H-D1b agree (V12 8 and D8).

\subsection*{Numbers}

{\scriptsize\raggedright
\begin{longtable}{p{0.317\linewidth}p{0.317\linewidth}p{0.317\linewidth}}
\toprule
Quantity & Custom simulator & ns-3 \\
\midrule
\endhead
D1 minus severity table, smallest cell & +0.021 & -0.013 \\
cells whose lower bound is below -0.01 & 0 of 32 & 8 of 32 \\
D1 minus severity table, cells with loss of 40\% or more (95\% interval) & +0.116 [+0.113, +0.120] & +0.021 [+0.017, +0.025] \\
D1 runs above the duty limit & 0 of 3200 & 0 of 3178 \\
D1 minus best baseline of the cell: smallest / mean & -0.060 / -0.026 & -0.070 / -0.005 \\
cells below the best baseline by more than 2 points (lower bound) & 26 of 32 & 12 of 32 \\
\bottomrule
\end{longtable}}

Mean Critical delivery within 10 s of the crash over the cells:

{\scriptsize\raggedright
\begin{longtable}{p{0.317\linewidth}p{0.317\linewidth}p{0.317\linewidth}}
\toprule
Policy & Custom & ns-3 \\
\midrule
\endhead
P1 (ns-3: p1) & 0.612 & 0.621 \\
P0 & 0.607 & not implemented in ns-3 \\
P4 & 0.637 & not implemented in ns-3 \\
quick & 0.713 & 0.510 \\
tamang & 0.712 & not implemented in ns-3 \\
khan & 0.592 & not implemented in ns-3 \\
nessa & 0.670 & 0.571 \\
D1 (ns-3: d1) & 0.703 & 0.631 \\
standard LoRaWAN confirmed uplink (NbTrans 8, RX1 1 s, RX2 2 s, ADR off); custom column: Khan-style & 0.592 & 0.557 \\
\bottomrule
\end{longtable}}

The Nessa-style row is ADDED-AFTER-SESSION-1 (V12 11.1; parameters frozen on the DEV grid of the custom simulator and reused in ns-3; implemented from a secondary description).

\textbf{Decision rule A4 (V12 11.4): D1 goes to the main paper only if H-D1b and H-D1s hold in ns-3 as well.} H-D1b in ns-3: not met; H-D1s in ns-3 (best of p1, quick, lorawan, nessa): not met. Result: \textbf{D1 is reported as simulator-dependent and stays in the supplement}.

\subsection*{What differs between the two simulators (stated, not isolated)}
\begin{itemize}
\item The custom policies use a 14-byte flash before every alert round, and a received flash counts as delivery; the ns-3 policies have no flash frame. The D1 table also plans one 97-byte transmission per attempt. Part of D1's gain in the custom simulator may come from the flash; this was not tested.
\item ns-3 ACK frames are about 13 bytes (the custom ACK is 26); receive windows are set to 0.1 s and 0.2 s for p1, quick and d1 (a local change) and kept at 1 s and 2 s for standard LoRaWAN.
\item The ns-3 baselines p1 and quick are simplified re-implementations from the descriptions. ADR on was not run (NOT\_RUN): in a one-hour run with one uplink per ten minutes ADR has no data to act on.
\item ns-3 allows one alert per vehicle at a time (a vehicle in its backoff tail gives way to a new crash), the custom simulator keeps several.
\item Different collision and reception models (ns-3's own interference matrix); different random draws.
\end{itemize}

\subsection*{Reading}

Where the verdicts differ, D1's gain over the severity table is simulator-dependent and is reported so. The comparison with published numbers (X1.4: ICC 2018 ns-3 paper, IEEE TWC 2021 single-gateway model, IEEE Sensors J. 2022 real-versus-simulated study) is NOT\_RUN: none of the three was opened. Absence statements: no claim.

